\section{Compressed generators and adaptive alphabets}
\label{rec:sec-adaptive}

A fixed representation level does not determine the dimension: diagrams
with the same first row can have very different shapes below it.  We now
separate these two contributions in another way.  Compression of the
transposition generator gives the level estimate.  A fourth trace moment
then gives a negative power of the full dimension, at a positive cost
proportional to the level scale.  Combining the two cancels that cost.

The fourth-moment argument realizes a diagram in an induced tensor
representation.  Its alphabet size controls the multiplicity paid at each
split.  An alphabet that was economical at the root can become too large
for a small descendant diagram.  We change it at that point and prove that
the total cost of all such changes is summable.  The proof below first
establishes the estimate for one split, then specifies the changes and
accounts for their cost over the entire recursion tree.

Put $F_n=T_nT_n^*$, so that
$F_n=P(F_m\otimes F_m)P$ for $n=2m$.  If
$\lambda=(n-r,\rho)\vdash n$, write
$h_n(r)=r\log(n/r)+r$ for $r>0$, and $h_n(0)=0$.
The coordinate names here are ordered so that $P$ is the final matching
of $T_n$.  Thus $F_n$ is the other positive orientation of the sweep,
with the same singular values as $T_n^*T_n$.  All traces are unnormalized;
$0\log0=0$ throughout.

\begin{theorem}\label{rec:adaptive-main}
There are absolute positive constants $c,c',C$ such that for every
dyadic $n$ and every partition $\lambda=(n-r,\rho)\vdash n$,
\begin{align}
 \|F_n|_{V_\lambda}\|_{\op}&\le e^{-c h_n(r)},
       \label{rec:generator-level-main}\\
 \Tr_{V_\lambda}F_n^4&\le
       \exp\{-c'\log D_\lambda+C h_n(r)\}.
       \label{rec:adaptive-trace-main}
\end{align}
Consequently $\|T_n|_{V_\lambda}\|_{\op}$ is at most
$\exp[-c_0\{\log D_\lambda+h_n(r)\}]$ for an absolute $c_0>0$.
\end{theorem}

\subsection{Compression of the transposition generator}

Let
\[
 L_n=\frac1n\sum_{i<j}(I-(ij)),\qquad
 T=L_n-L_0-L_1,\qquad J=L_0-L_1,
\]
where $L_0$ and $L_1$ act on the first and second halves, respectively,
and each uses normalization $1/m$ when $n=2m$.  Transpositions in these
formulas denote their unitary representation matrices.
The transposition content formula~\cite{VershikOkounkov2005} makes
$L_n$ scalar on $V_\lambda$, with value
\begin{equation}\label{rec:central-level}
 \ell=r-\frac{\binom r2+c(\rho)}n,
 \qquad r/2\le\ell\le r.
\end{equation}
Here $c(\rho)$ is the sum of the box contents of $\rho$.
The three operators in this display commute with each other before
compression, although they need not commute with $P$.

\begin{lemma}\label{rec:compressed-mgfs}
Let $1\le r\le n-1$ and let $P$ act on the ordered injective
$r$-tuple module by simultaneously switching matched positions.
Put $p_*=r/n$.  For absolute
$t_0,t_1,v_0>0$,
\begin{align}
 \|Pe^{tT}P\|&\le e^{Cp_*|t|},&&|t|\le t_0,
       \label{rec:T-small-mgf}\\
 \|Pe^{\pm t_1\log(1/p_*)T}P\|&\le1+o(1),&&p_*\longrightarrow0,
       \label{rec:T-log-mgf}\\
 \|Pe^{vJ^2/r}P\|&\le1+O(v),&&0\le v\le v_0.
       \label{rec:J-square-mgf}
\end{align}
The estimates are uniform in $n,r$.
\end{lemma}
\begin{proof}
Embed injections in the tensor power of the one-site space and let
$D$ be the mask deleting repeated sites.  Write $s_j=1$ or $-1$
according as the $j$th coordinate lies in the first or second half.
The one-coordinate part of $T$ is a projection onto the signed uniform
vector.  Its correction on distinct tuples is
$n^{-1}\sum_{j<l}s_js_l(I+(jl))$.
The one-coordinate part of $J$ is the sign diagonal minus the
difference of the two half-uniform rank-one projections; its correction
is $-(2m)^{-1}\sum_{j<l}(s_j+s_l)(I+(jl))$.

Use the symmetric and antisymmetric basis within matched site pairs.
A pair site is odd when an odd number of its tensor slots use the
antisymmetric vector.  The endpoints selected by $P$ have no odd pair
sites; two antisymmetric slots at the same pair site have even parity.
The injection mask kills basis indices with more than two tensor slots
at one pair site and preserves the odd-site pattern.  At a pair site
occupied by slots $j,l$, it is the projection
$(I-s_js_l)/2$ in the original position basis.  In the symmetric and
antisymmetric basis the second term flips both slot signs, so this
projection has absolute row sum one.  The product over disjoint doubly
occupied sites has the same bound.  Hence mask insertions between steps
have absolute row mass at most one.

Here is the transition count in this basis.  Write $e_{p,+},e_{p,-}$
for the two vectors at pair site $p$, and let $S_a$ flip the sign of
slot $a$.  Let $\tau_{ab}$ exchange the two tensor slots, so the
correction to $T$ is $n^{-1}\sum_{a<b}S_aS_b(I+\tau_{ab})$.
Only states with at most two slots per pair site survive the mask.
If exactly $k$ pair sites are odd, at most $2k$ slots lie at odd sites.
The mask preserves the entire set of odd sites, as well as the slot
occupancies, so the following bounds can be checked before each mask
insertion.

The one-slot rank-one part of $T$ sends a minus slot from a site $p$
to a minus slot at any site $q$, with coefficient $1/m$, and kills
plus slots.  When $p\ne q$, it toggles the parity at precisely those
two sites.  Even--even moves have total absolute mass at most
$rm/m=r$.  Odd--odd moves have at most $2k$ source slots and $k$
target sites, giving $2k^2/m$.  For moves preserving $k$, an odd
source and even target cost at most $2k$, whereas an even source and
odd target cost at most $rk/m$.  The same-site moves have total mass
at most $r/m$.  Thus this last term is $O(p_*)$, even when many
occupied sites have even parity.

The two-slot correction toggles the parities of the two sites holding
$a,b$; the slot exchange changes neither their occupancies nor these
parity changes.  Its two summands have total coefficient at most $2/n$
per slot pair.  Pairs at two even sites give mass $O(r^2/n)=O(r)$;
pairs at two odd sites give $O(k^2/n)$; and pairs at one site of each
parity give $O(kr/n)=O(k)$.  A pair of slots at the same site leaves
its parity unchanged.  There are at most $r/2$ such pairs, so their
mass is $O(r/n)=O(p_*)$.  Combining the two parts, and using
$r/m\le2$, gives the absolute transition bounds
\[
 \begin{array}{c|ccc}
 \text{change in odd-site count}&+2&-2&0\\ \hline
 T&Cr&Ck^2/m&C(k+p_*)
 \end{array}
\]
for both rows and columns by self-adjointness.  In particular the
fixed-count cost at $k=0$ is $O(p_*)$, not $O(r)$.
Weight level $k=2z$ by
$z!(rm)^{-z/2}$.  The off-diagonal masses then become at most
$C\sqrt{p_*}(z+1)$ upward and $C\sqrt{p_*}z$ downward.
Every return path with $2u$ off-diagonal steps has exponential-series
weight bounded by
\[
 e^{Cp_*|t|}
 \frac{(C\sqrt{p_*}|t|(u+1))^{2u}}{(2u)!}
 e^{C|t|u}.
\]
There are at most $4^u$ jump-direction lists with $2u$ up or down
steps; their intervening stays sum to the displayed exponential factor.
The weights telescope to one on every path returning to level zero.
Since $(2u)!\ge(2u/e)^{2u}$, the series excluding its zero-step term
is bounded by a geometric series with ratio
$C'p_*t^2e^{C|t|}$.  For fixed sufficiently small $|t|$ it is
$O(p_*t^2)$, which is absorbed in $e^{Cp_*|t|}$.
The matrices are self-adjoint, so the same absolute row bound applies to
columns and Schur's test proves~\eqref{rec:T-small-mgf}.
At $|t|=t_1\log(1/p_*)$, this ratio tends to zero if $t_1$ is a
sufficiently small positive constant, while $p_*|t|\to0$.
This proves~\eqref{rec:T-log-mgf} uniformly in $n,r$.

For $J$, the one-slot sign diagonal flips one sign at its current
site.  The difference of the half-uniform projections sends
$e_{p,+}$ to $m^{-1}\sum_qe_{q,-}$ and $e_{p,-}$ to
$m^{-1}\sum_qe_{q,+}$.  In the first case only the target site's
parity changes, and in the second only the source site's parity
changes.  Thus all these moves change $k$ by exactly one.  Their
upward mass is $O(r)$; downward moves either use one of the at most
$2k$ slots at odd sites or choose one of the $k$ odd target sites,
for total $O(k+rk/m)=O(k)$.
The correction $-n^{-1}\sum_{a<b}(S_a+S_b)(I+\tau_{ab})$
also toggles just one site's parity.  Its full mass is $O(r^2/n)$,
and its downward mass is $O(kr/n)$ because the toggled site must be
odd.  These are respectively $O(r)$ and $O(k)$, including coincident
pair sites.  Mask insertions preserve both conclusions.
Consequently $J$ has upward and downward masses at most $Cr$ and
$Ck$.  A return of length $2j$
therefore has total mass at most $(Crj)^j$.  Expanding
$e^{vJ^2/r}$ and using $j^j/j!\le e^j$ gives
\eqref{rec:J-square-mgf}.  The mask can be inserted between every
step because its absolute row mass does not exceed one.  Thus these
estimates apply to injections, not only to unrestricted tensors.
\end{proof}

\begin{proof}[Proof of \eqref{rec:generator-level-main}]
We first obtain a uniform estimate without the logarithmic factor.
Suppose at size $m$ that $F_m\preceq e^{-a_m L_m}$ on every
irreducible, with $0<a_m\le t_0$.  On the parent $\lambda$-space,
centrality and~\eqref{rec:T-small-mgf} imply
\[
 \|P(F_m\otimes F_m)P\|
 \le e^{-a_m\ell}\|Pe^{a_mT}P\|
 \le e^{-a_m\ell+Ca_mr/n}
 \le e^{-a_m(1-2C/n)\ell}.
\]
At one fixed sufficiently large dyadic size a positive $a_m$ is
available from the strict finite-size norm gap.  Starting above $4C$,
the product of the factors $1-2C/n$ over doubled sizes is positive,
since the sum of their deficits is finite.  Thus for an absolute
$a_0>0$ we have $\|F_n|_{V_\lambda}\|\le e^{-a_0\ell}$;
finitely many smaller sizes are absorbed by reducing $a_0$.

This estimate will pay an $O(r)$ error.  To obtain the additional
$\log(n/r)$ factor, we strengthen it by inducting on
\begin{equation}\label{rec:generator-penalty-induction}
 \|F_n|_{V_\lambda}\|
 \le\exp\{-aH_n(\ell)+D_0\max(0,2r-1)\},
 \qquad H_n(x)=x\log(n/x)+x.
\end{equation}
Take a unit vector in the range of $P$ in a minimal $r$-tuple
realization of $V_\lambda$.  Decompose it by the set of labels in each
half and by the two child position types.  Squared component norms give
the spectral probability law used below.  If the occupancies are
$j_0+j_1=r$ and the child first-row deficits are $r_0,r_1$, then
$0\le r_i\le j_i$, hence $r_0+r_1\le r$.  Let $\ell_0,\ell_1$
be the corresponding child generator values and put
$z_i=\ell_i/\ell$.  Since $\ell_i\le r_i\le r\le2\ell$, the
variables $z_i$ lie in $[0,2]$.  The decrement is
\[
 Z=H_n(\ell)-\sum_iH_m(\ell_i)
 =\left(\ell-\sum_i\ell_i\right)\log(n/\ell)
 +\ell\sum_i\{z_i\log(2z_i)-z_i+1/2\}.
\]
Since $z_i\in[0,2]$ and
$z\log(2z)-z+1/2\le C(z-1/2)^2$ there,
\begin{equation}\label{rec:generator-decrement}
 Z\le(\log(n/\ell)+C)|T|+CJ^2/r.
\end{equation}
For a unit vector in the range of $P$, spectral expectation in the
commuting child algebra and Lemma~\ref{rec:compressed-mgfs} give
\begin{equation}\label{rec:generator-slack}
 \E e^{2aZ}\le3/2\qquad(r/n\le p_0)
\end{equation}
for sufficiently small absolute $a,p_0$.  For the absolute value in
\eqref{rec:generator-decrement}, use
$(e^{|u|/2}-1)^2\le C(\cosh u-1)$ and both signs in
\eqref{rec:T-log-mgf}; then use Cauchy--Schwarz and
\eqref{rec:J-square-mgf}.

For $r\ge2$, the sum of the child penalties in
\eqref{rec:generator-penalty-induction} drops by at least one,
except when $(r_0,r_1)$ is $(r,0)$ or $(0,r)$.  This persistence
requires $j_0=r$ or $j_1=r$, respectively.  On the range
of $P$, this persistence event has weight at most $2^{1-r}$.
Indeed it requires all tuple coordinates in one half.  If two occupy
a matched pair it is impossible; otherwise independent switches give
that probability exactly.  The ratio in the induction is consequently
at most
\[
 \sqrt{(3/2)2^{1-r}}+e^{-D_0}\sqrt{3/2}\le1
\]
when $D_0$ is large.  Level zero is trivial and level one is annihilated because a single
card is uniform after one sweep.  In the range $r/n>p_0$, the ratio
$H_n(\ell)/r$ is bounded by a constant depending only on $p_0$,
so a sufficiently large $D_0$ makes the right side of
\eqref{rec:generator-penalty-induction} at least one.  This starts
and closes that induction, with bounded sizes covered by the same
finite choice of constants.

It remains to remove the positive penalty.  For a fixed sufficiently
small $\theta>0$, geometrically interpolate the bound
$e^{-a_0\ell}$ with~\eqref{rec:generator-penalty-induction}.
Because $2D_0r\le4D_0\ell$, choose $\theta$ so that
$4\theta D_0\le(1-\theta)a_0$.  The penalty is then canceled and
we obtain $\|F_n|_{V_\lambda}\|\le e^{-\theta aH_n(\ell)}$.
Finally $r/2\le\ell\le r$ gives
$H_n(\ell)\ge\tfrac12h_n(r)$ and proves the level estimate.
\end{proof}

\subsection{The induced tensor spaces and a single split}

The level estimate is now complete.  The remaining goal is the
fourth-moment bound.  We will prove it on larger representations containing
$V_\lambda$, because a trace of a positive power can only increase under
such an enlargement.

More precisely, a list of categories with sizes $b_j$ and representations
$V_{\eta_j}$ defines
\[
 \mathcal M=\operatorname{Ind}_{\prod_jS_{b_j}}^{S_n}
                    \bigotimes_jV_{\eta_j},\qquad \sum_jb_j=n.
\]
Its log trace is $\sigma=\log\operatorname{Tr}_{\mathcal M}F_n^4$,
with $\sigma=-\infty$ if the trace is zero.  Empty categories may be
omitted.  An induced module will have categories of total size $n$.  A row
category consists of a diagram $\eta$ of size $b$, an alphabet of
dimension $q\ge\ell(\eta)$, where $\ell(\eta)$ is its number
of rows, and an orientation, ordinary or
sign-twisted.  A regular category consists of a spatial diagram $\kappa$ of
size $b$ and $b$ distinct named labels.  We may designate one ordinary
one-row category of alphabet dimension one as the \emph{background};
it may be empty.  All other categories are \emph{payload}, including
ordinary one-row categories of alphabet dimension one.  In the
application the background is the original first row of $\lambda$.
Every descendant of a payload category remains payload after
restriction or a change of alphabet.  There is at most one background
category in each node, so its child counts are determined by the payload
counts and the two half sizes.  Isomorphic categories remain distinct
when their labels or origins differ.

For a row category $j$, $\eta$ denotes the alphabet type before the
sign twist.  Its spatial factor $V_{\eta_j}$ in $\mathcal M$ is
$V_\eta$ in ordinary orientation and
$V_\eta\otimes\operatorname{sgn}_{S_b}\cong V_{\eta'}$ in sign-twisted
orientation, where $\eta'$ is the conjugate partition.  In both
orientations, $q\ge\ell(\eta)$ and the multiplicity $G_q(\eta)$ below
refer to this pre-twist diagram.

Realize row categories in signed tensor powers.  Each category has its
own alphabet; let $\mathcal H$ be the space spanned by words specifying
one alphabet symbol at each of the $n$ position slots, with the prescribed
number of slots in each category.
An adjacent swap contributes $-1$ exactly when both symbols are odd;
these adjacent swaps generate the graded permutation action.  Use odd symbols precisely for sign-twisted categories.  This gives
the indicated representation on the stabilizer of a category allocation,
and induction sums over allocations.  In a regular category, each named
label appears once.  The commuting group is $U(q)$ for each row alphabet
and $S_b$ acting on the names of each regular category.  Let $E$ be
the product of their isotypic projections onto the row types $\eta$
and regular types $\kappa$.  On a fixed category allocation,
Schur--Weyl duality and the regular bimodule decomposition give the
spatial tensor product and its multiplicity space.  Summing over
allocations therefore gives
\begin{equation}\label{rec:alphabet-multiplicity}
 g=\prod_{\mathrm{row}}G_q(\eta)
                    \prod_{\mathrm{regular}}D_\kappa
\end{equation}
copies of the desired spatial induced module:
\[
 E\mathcal H\cong\mathbb C^g\otimes\mathcal M.
\]
Here $G_q$ is the dimension of the polynomial $U(q)$ representation.
Spatial permutations act trivially on $\mathbb C^g$, so $E$ commutes
with the network operators and a spatial trace on $E\mathcal H$ is
$g$ times the trace on $\mathcal M$.

A child projection $W$ specifies counts $b_0+b_1=b$ and types
$\mu,\nu$ in the two halves of each category.  For a row category,
use the two child $U(q)$ isotypic projections, which commute with
the diagonal $U(q)$ action and hence with $E$.  For a regular
category first fix its subset of $b_0$ names in the first half and
project the two name groups onto $\mu,\nu$; then sum these orthogonal
projections over all $\binom b{b_0}$ subsets.  This sum preserves each
name subset and is invariant under global relabeling, so it also
commutes with $E$.  The projections $W$ sum to the identity and
commute with $Q=F_m\otimes F_m$.  Only types with positive
Littlewood--Richardson coefficient for the parent type contribute
to $WE$; we call these choices compatible.  All row types in this
description are taken before their inherited sign twist.  Write $g_i$
for the multiplicity \eqref{rec:alphabet-multiplicity} evaluated on
the induced module in half $i$.
Put
\[
 \mathcal L=n\log n-
 \sum_{\mathrm{row}}\sum_i\eta_i\log\eta_i
 -\sum_{\mathrm{regular}}b\log b.
\]

\begin{lemma}\label{rec:alphabet-angle}
For compatible parent and child types,
\[
 \|WPE\|^2\le
 \prod_{\mathrm{row}}G_q(\mu)G_q(\nu)\,
       e^{-(\mathcal L-\mathcal L_0-\mathcal L_1)/2}.
\]
\end{lemma}
\begin{proof}
Choose block density matrices $A,B$ on the one-site alphabet.  On a
row block their spectra are $(\mu_i/m)$ and $(\nu_i/m)$; on a
regular block they are scalar $b_0/(mb)$ and $b_1/(mb)$.
Their traces, summed over all blocks, are one.  If $V$ applies
$A^{1/2}$ in the first half and $B^{1/2}$ in the second, then
\begin{equation}\label{rec:alphabet-geometric-mean}
 PVP\le P\left(((A+B)/2)^{1/2}\right)^{\otimes n}P.
\end{equation}
This is the needed joint-concavity instance for the operator geometric
mean; see also~\cite[Appendix A.2.1, Lemma~4]{FawziSaunderson2020}.
We include its block-matrix proof.  On one matched pair, average
the two positive block matrices with diagonal entries $A\otimes I,I\otimes B$ and
$B\otimes I,I\otimes A$.  Their positive off-diagonal blocks are
the respective geometric means.  The maximal positive off-diagonal
block with diagonal entries $X,Y$ is
$X^{1/2}(X^{-1/2}YX^{-1/2})^{1/2}X^{1/2}$, by the Schur complement
and monotonicity of square roots.  Averaging and then tensoring proves
\eqref{rec:alphabet-geometric-mean}; continuity handles zero eigenvalues.
The signed flips cause no change because the matrices respect category
blocks.

For one row category, average the eigenbasis of $A$ over $U(q)$.
On child type $\mu$, Schur's lemma makes the averaged operator
$(A^{1/2})^{\otimes b_0}$ scalar on its multiplicity space.  Its scalar is
the Schur character $s_\mu$ divided by $G_q(\mu)$, and the highest-weight
monomial gives the lower bound
\[
 \frac{s_\mu(\sqrt{\mu_1/m},\ldots,\sqrt{\mu_q/m})}{G_q(\mu)}
 \ge \frac{\prod_i(\mu_i/m)^{\mu_i/2}}{G_q(\mu)}.
\]
Zero exponents contribute one.  Apply the same calculation to $B$
and $\nu$, independently in every row block and in the two halves.
A regular block contributes the scalar
$(b_0/(mb))^{b_0/2}(b_1/(mb))^{b_1/2}$.  Multiplication therefore
shows that the average of $V$ dominates $W$ times
\[
 \frac{\exp[-(\mathcal L_0+\mathcal L_1+
                      \sum_{\mathrm{regular}}b\log b)/2]}
      {\prod_{\mathrm{row}}G_q(\mu)G_q(\nu)}.
\]
For the parent bound put $C=(A+B)/2$.  If $x_{j,i}$ are its
eigenvalues in nonincreasing order in each row block, the norm of
$(C^{1/2})^{\otimes b}$ on
parent type $\eta$ is the highest-weight product
$\prod_i x_{j,i}^{\eta_i/2}$.  The regular blocks of $C$ are each
$1/n$ times the identity on their $b$ names.  All eigenvalues of
$C$ sum to one.  The entropy maximization
\[
 \prod_{\mathrm{row},i}x_{j,i}^{\eta_{j,i}/2}\,n^{-s/2}
 \le \prod_{\mathrm{row},i}(\eta_{j,i}/n)^{\eta_{j,i}/2}
                  n^{-s/2}
\]
follows by maximizing subject to the remaining row trace $1-s/n$,
where $s$ is the regular mass.  Thus on $E$ the right side of
\eqref{rec:alphabet-geometric-mean} has norm at most
$\exp[-(\mathcal L+\sum_{\mathrm{regular}}b\log b)/2]$.
This upper bound is uniform in the eigenbases just averaged.  Compress
the averaged inequality by $E$ and use
$EPWPE=(WPE)^*(WPE)$.  Dividing the parent bound by the child scalar
proves the lemma.  Sign twists change only the spatial action, so the
same multiplicity-space calculation applies in either orientation.
\end{proof}

Take $p=4$, and let $Y=WQ\ge0$.  The needed Schatten inequality is
\begin{equation}\label{rec:alphabet-strip}
 \Tr(EPYPE)^p
 \le\|WPE\|^{2p-4}\Tr(PY^{p/2}P)^2.
\end{equation}
Indeed apply three-lines interpolation to $Y^{pz/4}WPE$, extended
by zero on the kernel of $Y$, between the operator norm at $\Re z=0$
and the fourth Schatten norm at $\Re z=1$.  At $z=2/p$ the target
norm is the $2p$ norm.  Raising the interpolation inequality to $2p$
gives \eqref{rec:alphabet-strip}: at the fourth-norm endpoint first
obtain $\Tr(EPY^{p/2}PE)^2$, then drop the orthogonal compression
by $E$ using Hilbert--Schmidt contraction.  Thus no commutation of $P$
with $Y$ is required.

\begin{lemma}\label{rec:alphabet-cycle-inflation}
Let $s$ be the total number of named labels and
$\sigma_i=\log\Tr F_m^4$ on the child induced modules.  If $m\ge2$,
\begin{equation}\label{rec:alphabet-pair-trace}
 \Tr(PY^2P)^2\le
 g_0g_1e^{\sigma_0+\sigma_1}\,
 2^{\gamma s/2}\prod_{\mathrm{regular}}
                  \left(2^{-b}\binom b{b_0}\right),
 \qquad\gamma=\frac{\log_2 3}{2}<1.
\end{equation}
For $m=1$ the inflation $2^{\gamma s/2}$ may be replaced by $2^{s/2}$.
\end{lemma}
\begin{proof}
Collapse each matched pair of tensor slots to its unordered contents.
On such an orbit, $P$ has rank at most one and its nonzero vector has
constant absolute coefficient.  For orbit states $C,D$, let $j(C,D)$
be the number of cycles in the union of the two partial matchings on
named labels, counting a common edge twice as a cycle.  The fraction
of orientation pairs agreeing on the labels' half assignments is
$2^{-s+j(C,D)}$.  Indeed, if the two matchings have $a,a'$ edges,
their separate half assignments number $2^{s-a}$ and $2^{s-a'}$.
The union has $2^{s-a-a'+j(C,D)}$ assignments, since each cycle
has one redundant opposite-half condition.  Dividing gives the
fraction.  Cauchy--Schwarz
therefore bounds the squared orbit entry by
\begin{equation}\label{rec:alphabet-orbit-entry}
 2^{-s+j(C,D)}\sum_{u,v}|Y^2(u,v)|^2.
\end{equation}

Here the sum is over word pairs in the two orbits; entries of $Y^2$
vanish unless each named label stays in its assigned half.  For such
a word pair $u,v$, let $A_*,B_*\subseteq[m]$ be the first-half sites
occupied by names in $u,v$, respectively, and let $C_*,D_*$ be the
corresponding second-half site sets.  The partial bijections
$f:A_*\to B_*$ and $g:C_*\to D_*$ follow each named label from
$u$ to $v$.  A subset $U\subset A_*\cap C_*$ satisfies $f(U)=g(U)$
exactly when it is a union of closed components of the resulting
partial permutation.  These components are the cycles counted by
$j(C,D)$, so the number of such subsets is $2^{j(C,D)}$.

Let $\tau$ range uniformly over the last matching switches in the
second child, and apply it to the second-half sites of the output
word $v$.  This replaces $g$ by $\tau g$, leaving $f$ fixed.
The squared kernel weight is unchanged: $W$ commutes with these
spatial permutations, and the positive power of $Q$ has range fixed
by them.  For each $U$, the probability of
$f(U)=\tau g(U)$ is either zero or
$|\mathcal O(g(U))|^{-1}$, where $\mathcal O$ is the matching
group orbit.  Thus averaging inside the sum over word pairs bounds
the subset count by
\[
 \sum_{X\subset g(A_*\cap C_*)}|\mathcal O(X)|^{-1}.
\]
A switch pair with two available slots contributes
$1+2(1/2)+1=3$; one available slot contributes $1+1/2\le\sqrt3$.
Since $|A_*\cap C_*|\le\min(s_0,s_1)\le s/2$, where $s_i$ is
the number of names in half $i$, the product is at most
$3^{s/4}=2^{\gamma s/2}$.  Sum
\eqref{rec:alphabet-orbit-entry} over words.  For a fixed choice of
name subsets the unswitched trace is $g_0g_1e^{\sigma_0+\sigma_1}$.
The aggregated projection is their orthogonal direct sum, giving
\begin{equation}\label{rec:alphabet-aggregated-trace}
 \Tr Y^4=\Tr(WQ^4)
 =\left(\prod_{\mathrm{regular}}\binom b{b_0}\right)
                 g_0g_1e^{\sigma_0+\sigma_1}.
\end{equation}
Together with the factor $2^{-s}$ this proves the result.  The named
subsets enter through this trace identity exactly once.
At $m=1$, simply use $j(C,D)\le s/2$.
\end{proof}

We can now pass from the count-space trace to the desired spatial
trace.  Let $M$ be the number of compatible child projections $W$.
For each payload category of size $b$, the count split and two child
partitions have at most $\exp(Cb^{2/3})$ choices, by the partition
bound.  The single background has no type choice and its counts are
determined by the payload counts.  Therefore
$\log M\le C\sum_{\mathrm{payload}}b^{2/3}$.  Named subsets
are already aggregated in $W$ and are not choices counted by $M$.

Resolve $Q=\sum_W WQ$ on $E\mathcal H$.  The Schatten triangle
inequality, followed by \eqref{rec:alphabet-strip} with $p=4$,
the angle bound, and \eqref{rec:alphabet-pair-trace}, gives, for $m\ge2$,
\begin{align}
 g e^\sigma
 &=\Tr(EPQPE)^4\notag\\
 &\le M^4\max_W\bigg\{
 e^{\sigma_0+\sigma_1-(\mathcal L-\mathcal L_0-\mathcal L_1)}
 \left[\prod_{\mathrm{row}}G_q(\mu)G_q(\nu)\right]^2
 \notag\\[-2pt]
 &\hspace{5em}{}\times g_0g_1 2^{\gamma s/2}
 \prod_{\mathrm{regular}}2^{-b}\binom b{b_0}\bigg\}.
 \label{rec:alphabet-master-split}
\end{align}
The angle contributes two copies of each child row carrier, and the
child trace contributes a third.  Dividing by $g$ and taking logarithms
therefore yields the explicit bound
\begin{align*}
 \sigma\le\max_W\bigg\{&\sigma_0+\sigma_1+4\log M
      -(\mathcal L-\mathcal L_0-\mathcal L_1)\\
 &+\sum_{\mathrm{row}}
       [3\log G_q(\mu)+3\log G_q(\nu)-\log G_q(\eta)]\\
 &+\sum_{\mathrm{regular}}
       [\log D_\mu+\log D_\nu-\log D_\kappa
          +(\gamma/2-1)b\log2+\log\tbinom b{b_0}]
 \bigg\}.
\end{align*}
For $m=1$ replace $\gamma$ by one, at an extra cost $O(s)$.
This calculation motivates the potentials that we now telescope.

For a diagram of size $b$, put
$I(\eta)=\sum_i\eta_i\log(b/\eta_i)$ and set $\delta=1/32$.
Give categories the potentials
\begin{equation}\label{rec:alphabet-potentials}
 \phi(\eta)=\delta I(\eta)\quad\text{(row)},\qquad
 \phi(\kappa)=\log D_\kappa-\frac{1-\delta}{2}b\log b
                                     \quad\text{(regular)}.
\end{equation}
Let $\Phi$ be their sum.  For every category with child counts $b_0+b_1=b$, define
\[
 v_b=b\{\log2-H_2(b_0/b)\},\qquad
 H_2(x)=-x\log x-(1-x)\log(1-x),
\]
with $v_0=0$.  Thus $H_2$ uses natural logarithms in this section,
and $v_b$ is the loss from an imbalanced count split.
The preceding lemmas imply the single-split accounting inequality
\begin{align}
 \sigma\le\max\bigg\{&\sigma_0+\sigma_1
 -(\Phi-\Phi_0-\Phi_1)-\sum_{\mathrm{payload\ row}}v_b
 \notag\\
 &+C_1\sum_{\mathrm{payload}}b^{2/3}
 +C_1\sum_{\substack{\mathrm{payload\ row}\\q\ge2}}
       q^2\log(e(1+b/q^2))+C_1s\mathbf1_{m=1}\bigg\}.
       \label{rec:alphabet-split-budget}
\end{align}
Here the maximum is over compatible child counts and types.
To derive this form from \eqref{rec:alphabet-master-split}, use the
Cauchy identity
\[
 \sum_{\substack{\xi\vdash b\\\ell(\xi)\le q}}G_q(\xi)^2
       =\binom{b+q^2-1}{q^2-1}.
\]
It bounds the positive child row-carrier logarithms by
$Cq^2\log(e(1+b/q^2))$; for $q=1$ all carrier factors are one.
Discarding the favorable parent carrier logarithm only increases
the bound.  Also
\[
 \mathcal L-\mathcal L_0-\mathcal L_1
 =\sum_{\mathrm{all}}v_b+
   \sum_{\mathrm{row}}(I(\eta)-I(\mu)-I(\nu)),
\]
and both sums have nonnegative summands.  Indeed compatibility implies
that $\mu+\nu$ dominates $\eta$; Schur concavity followed by entropy
concavity gives $I(\eta)\ge I(\mu+\nu)\ge I(\mu)+I(\nu)$.
Retain $-v_b$ for payload rows and replace the negative row-entropy
decrement by $-\delta[I(\eta)-I(\mu)-I(\nu)]$.  The unused
background and regular imbalances are nonnegative and may be discarded.
For a regular category the carrier ratio gives
$\log D_\mu+\log D_\nu-\log D_\kappa$, and
$\log\binom b{b_0}\le bH_2(b_0/b)$ gives
$(\gamma/2-1)b\log2+bH_2(b_0/b)
\le(1-\delta)bH_2(b_0/b)/2$ because $\delta<1-\gamma$.
This proves \eqref{rec:alphabet-split-budget}.

\subsection{Changing alphabets and paying for the changes}

The case $r=0$ is the trivial representation and satisfies both bounds
with equality.  Assume $r>0$.  Embed $V_\lambda$ in the module induced from its first row as
background and $\rho$ as payload.  Fix large absolute $A$ and then
a much larger $K$.  To tune a diagram of size $b$, take its rows
longer than $K\sqrt b$ as one row category; in the remaining diagram
take columns longer than $K\sqrt b$ as another row category, with
orientation reversed.  Put the residual diagram $S$ in a regular category.  Here is
the representation containment used in this operation.  If $R$ is the
selected top-row diagram and $\xi$ the remainder, the parent is their
sorted-row union, whose Littlewood--Richardson coefficient in
$\operatorname{Ind}(V_R\otimes V_\xi)$ is one.  Apply the transposed
version of this fact to remove the selected columns from $\xi$.
Transitivity of induction then embeds the parent in the module of
the three resulting parts.  Selected rows inherit the old orientation,
and selected columns reverse it.  If the old orientation was twisted,
the residual regular spatial diagram is $S'$ rather than $S$;
$D_{S'}=D_S$, so its potential is unchanged.  These conventions pass
the old sign twist to every spatial factor.  Thus tuning enlarges the
spatial module, and its positive-power trace can only increase.
Choosing a new alphabet afterward changes the multiplicity space in
which that spatial module is realized.
At the root only, group the selected row lengths $l$ by
$\lfloor\log_2(r/l)\rfloor$, keeping the two orientations separate.
No later tuning uses this grouping.

Assign a new row category its actual oriented height $q$.  A regular
category stays regular.  A row category of $q=1$ never changes.
For $q\ge2$, retain its alphabet while $b\ge Aq^2$; when
$0<b<Aq^2$, end that phase and tune the current diagram afresh.
At birth,
\begin{equation}\label{rec:alphabet-birth-exit}
 q_{\rm new}\le\sqrt b/K,\qquad
 b_{\rm new}\ge K^2q_{\rm new}^2;\qquad
 q_{\rm new}\le q/L\ \text{at an exit},\quad L=K/\sqrt A>1.
\end{equation}
Thus no immediate retuning is needed and alphabet dimensions shrink
geometrically along successive phases.

Let $b_j$ be the initial payload category sizes, $I_0$ the sum of their row
entropies, $s_0$ the initial regular size, and
$J_{\rm in}=\sum b_j\log(r/b_j)$.  Hook products and the long-row
construction give
\begin{align}
 \Phi_{\rm in}&\ge\delta I_0+\tfrac\delta2s_0\log s_0-C_Kr,
       \label{rec:alphabet-initial-potential}\\
 \log D_\lambda&\le\log\binom nr+I_0+\tfrac12s_0\log s_0+J_{\rm in},
       \label{rec:alphabet-initial-dimension}\\
 \sum_{j\ \mathrm{row}}b_j\log q_j&\le I_0+r\log2,\qquad
 J_{\rm in}\le\epsilon I_0+C_\epsilon r.
       \label{rec:alphabet-initial-entropy}
\end{align}
Here are the dimension and entropy estimates underlying these bounds.
The residual diagram has every row and column of length at most
$K\sqrt r$, so its hooks are at most $2K\sqrt r$ and
\[
 D_S\ge s_0!/(2K\sqrt r)^{s_0}.
\]
Together with $s_0\log(r/s_0)\le Cr$, this proves the first line.
For the dimension upper bound, induced-module containment contributes
$\binom nr$ and the payload allocation multinomial.  Bound these by
their entropy expressions, each row-category dimension by
$e^{I(\eta)}$, and the residual dimension by $\sqrt{s_0!}$.
This gives~\eqref{rec:alphabet-initial-dimension}.

Within each initial row category the positive row lengths differ by at
most a factor two.  Its row distribution therefore has entropy at least
$\log q_j-\log2$, proving the first assertion of
\eqref{rec:alphabet-initial-entropy}.  To prove the second, choose a
uniform payload box.  Its category is determined by a flag for the row,
column, or residual part and, for a row or column part, a nonnegative
dyadic group index $Z$, taking $Z=0$ on the residual part.  These
data have entropy $J_{\rm in}/r$.
The mean of $Z$ is at most
$(I_0+J_{\rm in})/(r\log2)$, because a row of length $l$ contributes
$\log_2(r/l)$ before taking the integer part.  For any $\alpha>0$,
relative entropy against the geometric law proportional to
$e^{-\alpha z}$ gives $H(Z)\le\alpha\mathbb EZ+C_\alpha$.
Adding the flag entropy, which is at most $\log3$, and choosing
$\alpha$ sufficiently small gives
$J_{\rm in}\le\epsilon I_0+C_\epsilon r$ after rearrangement.
This is the reason for grouping at the root; no subsequent tuning needs
that extra grouping.

We detail the amortization, because using one fixed alphabet would
leave a divergent carrier cost.  At an exit of mass $b$, the adverse
potential jump is at most
\begin{equation}\label{rec:alphabet-exit-jump}
 C_2b+C_Ks_{\rm new}+C_2b\log_+(q^2/b).
\end{equation}
To verify this jump estimate, denote by $R,C,S$ the new row, column,
and residual parts, of respective sizes $|R|,c_*,s$.  Entropy of the
part assignment gives
\[
 I_{\rm row}(\eta)\le I_{\rm row}(R)+I_{\rm row}(C)
                         +I_{\rm row}(S)+b\log3.
\]
Since the old height is at most $q$,
\[
 I_{\rm row}(C)-I_{\rm col}(C)
 \le c_*\log q-c_*\log(c_*/q),\qquad
 I_{\rm row}(S)\le s\log q.
\]
The hook bound for the residual gives
\[
 \phi(S)\ge\frac\delta2s\log b-C_Ks
                  -\frac{1+\delta}{2}s\log(b/s).
\]
Substituting these inequalities and using
$x\log(b/x)\le Cb$ for $0<x\le b$ proves
\eqref{rec:alphabet-exit-jump}; empty parts contribute zero.
The total mass ever converted to regular categories is at most $r$,
because a regular category never becomes a row category again.  We
therefore charge the $C_Ks_{\rm new}$ term only once to each converted
box, for a total at most $C_Kr$.  The linear term $C_2b$ in
\eqref{rec:alphabet-exit-jump} will be charged to the total exit mass.
Only its logarithmic excess needs an imbalance saving.

Call a split balanced when both child masses are in $[b/4,3b/4]$.
During a phase with $q\ge2$ the parent mass satisfies $b\ge Aq^2$.
At an unbalanced split,
$v_b\ge\theta b$, where $\theta=\log2-H_2(1/4)>0$.  Put
\[
 f(t)=\frac{\log(e(1+t))}{t},\qquad t>0.
\]
The carrier cost is $C_1bf(b/q^2)\le C_1bf(A)$, since $f$ is
decreasing.  If a child of mass $x$ exits at this split, then
\[
 x\log_+(q^2/x)\le q^2/e\le b/(eA).
\]
There are at most two exiting children.  Choose $A\ge4$ so large
that $C_1f(A)+2C_2/(eA)\le\theta/2$.  The carrier cost and these
logarithmic excesses are then bounded by half of the negative
imbalance term.  At a balanced exit, $x\ge b/4\ge Aq^2/4\ge q^2$,
so the logarithmic excess is zero.  No $K$-dependent constant has
entered this choice of $A$.

It remains to sum carrier costs at balanced nodes of a phase, which
may form a branching subtree.  Distribute each such cost evenly among
its $b$ boxes.  A box pays $C_1f(b/q^2)$.  Along its balanced ancestors,
successive masses decrease by a factor at most $3/4$, even if
unbalanced nodes intervene.  Reversing the sequence and using
$b/q^2\ge A$ gives
\[
 \sum_{\text{balanced ancestors}}f(b/q^2)
 \le\sum_{j\ge0}f(A(4/3)^j)
 \le \frac{C\log(e(1+A))}{A}.
\]
Summing over boxes bounds the balanced cost by a constant times the
phase's birth mass.  The constant depends on the fixed $A$ and is
independent of $K$.  The exit frontier has total mass at most that
birth mass, so it also pays the linear exit cost $C_2b$.
By \eqref{rec:alphabet-birth-exit}, a box initially in row category
$j$ encounters at most $1+\log q_j/\log L$ phases: each retuning
reduces the new alphabet by a factor $L$, and regular conversion ends
the sequence.  Each of the total birth mass and the total exit mass
is consequently at most
\[
 \sum_{j\ \mathrm{initial\ payload\ row}}b_j
             (1+\log q_j/\log L).
\]

The $b^{2/3}$ term is also summable, including categories that stay
small through many ambient levels.  At a balanced split,
\[
 b_0^{2/3}+b_1^{2/3}-b^{2/3}\ge c b^{2/3}.
\]
The power sum cannot decrease at any other split or conversion and its
terminal value is $r$, so all balanced costs sum to $O(r)$.
For unbalanced splits, the positive error satisfies the numerical bound
\[
 \sum_{\rm unbalanced}b^{2/3}
 \le\sum_{\rm unbalanced}b
 \le\theta^{-1}\sum_{\rm unbalanced}v_b.
\]
To bound the last sum, the scalar count entropy $b\log(n/b)$ at
ambient size $n$ decreases by exactly $v_b$ at a split.  A conversion
increases it by $\sum_j b_j^{\rm new}\log(b/b_j^{\rm new})$.
Its terminal value is zero, so summing all the nonnegative decrements
gives
\begin{equation}\label{rec:alphabet-count-telescope}
 \sum_{\mathrm{payload\ splits}}v_b
 =r\log(n_{\rm root}/r)+J_{\rm in}
 +\sum_{\rm exits}\sum_j b_j^{\rm new}\log(b/b_j^{\rm new}).
\end{equation}
This identity bounds the total positive error by the initial count
entropy and the conversion entropies; it makes no further subtraction
of the operator term $-v_b$.  Each exit creates at most three categories,
so its last sum is at most $b\log3$ and is covered by the phase-mass
bound.  Finally the $m=1$
term costs at most $Cr$.

Iterating \eqref{rec:alphabet-split-budget} down to single sites,
whose log trace is zero, now gives
\begin{equation}\label{rec:alphabet-telescoped}
 \sigma_{\rm in}\le-\Phi_{\rm in}+C_Kr+
 C_3\{r\log(n/r)+J_{\rm in}+r\}
 +\frac{C_3}{\log L}\sum_{j\ \mathrm{row}}b_j\log q_j.
\end{equation}
Here $C_3$ depends on the already fixed $A$, but not on $K$; all
$K$-dependent losses are in $C_Kr$.  Choose $K$ after $A$ so that
$C_3/\log L$ is small compared with
$\delta$, and then choose $\epsilon$ in
\eqref{rec:alphabet-initial-entropy} sufficiently small.
Equations \eqref{rec:alphabet-initial-potential}--
\eqref{rec:alphabet-telescoped} yield
\[
 \sigma_{\rm in}\le
 -\tfrac\delta2 I_0-\tfrac\delta2s_0\log s_0+C h_n(r).
\]
The dimension comparison \eqref{rec:alphabet-initial-dimension}
proves \eqref{rec:adaptive-trace-main}.  Combining its fourth-root
bound with \eqref{rec:generator-level-main}, in proportions chosen
to cancel the positive $C h_n(r)$ term, proves the final assertion
of Theorem~\ref{rec:adaptive-main}.

For clarity, this also supplies a complete mixing consequence.  Choose an
absolute integer $q$ so large that the dimension powers are absorbed in
\[
 \operatorname{Tr}_{\rm reg}F_n^q-1
 \le\sum_{r=1}^{n-1}2^r e^{-c_2q h_n(r)}=o(1),
\]
where $c_2>0$ is absolute after increasing $q$.  The count $2^r$
bounds partitions of the lower diagram.  Split at $r=\sqrt n$ to see
that the sum tends to zero: below it $h_n(r)\ge\frac12r\log n$,
and above it $h_n(r)\ge r$.  Schatten H\"older bounds the squared
Hilbert--Schmidt norm of the nonconstant part of $T_n^q$ by this
positive moment, without requiring normality.  Plancherel and translation
therefore give total-variation convergence from every initial deck after
$q$ sweeps, or $qd$ physical shuffles.
